% --- Archon physics formalization source begin ---
% archon:physics
% archon:covers PhyXMiniProblems/problem_phyx_mini_0147.lean
% archon:source-report reports/phyx_mini/problem_phyx_mini_0147.source.json

\chapter{Physics problem phyx\_mini\_0147}
\label{ch:PhyXMiniProblems_problem_phyx_mini_0147}

\paragraph{Problem source.}
\#\# Physical scenario

Figure shows a liquid-detecting prism device that might be used inside a washing machine. If no liquid covers the prism's hypotenuse, total internal reflection of the beam from the light source produces a large signal in the light sensor. If liquid covers the hypotenuse, some light escapes from the prism into the liquid and the light sensor's signal decreases. Thus a large signal from the light sensor indicates the absence of liquid in the reservoir.

\#\# Question

Determine the allowable range for the prism's index of refraction \textbackslash{}(n\textbackslash{}).

\#\# Answer choices

- A: A: \$1.41 < n < 1.98\$

- B: B: \$1.31 < n < 1.88\$

- C: C: \$1.41 < n < 1.88\$

- D: D: \$1.51 < n < 1.78\$

\#\# Figure caption (auxiliary; use the image as primary evidence)

This diagram depicts an optical setup involving two lenses and an object:

1. **Object**: On the left side, there is an upright arrow, representing an object. It is positioned on the horizontal axis with a label indicating it is 25.0 cm from the first lens.

2. **First Lens**: The first lens encountered by the light is a biconcave (diverging) lens. It is positioned along the horizontal axis and spaced 25.0 cm to the right of the object. 

3. **Second Lens**: To the right of the first lens, there is a biconvex (converging) lens. It is placed 30.0 cm from the diverging lens.

4. **Distances**: The distance from the second lens to the resulting image, represented by a dashed arrow pointing downward, is 17.0 cm.

5. **Light Ray Path**: The setup illustrates the path of light as it travels from the object through the lenses. The dashed arrow suggests the location of the image formed by the system, indicating that the image is inverted relative to the object.

The diagram includes numerical measurements to define the distances between the object, lenses, and image.

\#\# Dataset metadata

Geometrical Optics; Physical Model Grounding Reasoning, Multi-Formula Reasoning

\paragraph{Recorded answer/context.}
C: C: \$1.41 < n < 1.88\$

\paragraph{Figure/image path.}
/root/proposal\_for\_physic/science-mango/phyx\_mini\_run/phyx\_data/test\_image/147.png

\paragraph{Formalization target.}
create a compiling Lean file with sorry bodies at `PhyXMiniProblems/problem\_phyx\_mini\_0147.lean`. The Lean declarations must preserve the physical quantities, dimensions or dimensional roles, figure labels, governing-law hypotheses, and final relation expressed by this problem.
Use Mathlib/Physlib names found through LeanExplore where available. If a physics API is missing, introduce faithful local abstractions rather than scalar placeholder aliases.

\begin{theorem}[Physics formalization target]
\label{thm:physics:phyx_mini_0147:target}
\lean{PhyXMiniProblems.ProblemPhyXMini0147.problem_phyx_mini_0147}
\uses{def:physics:phyx-mini-0147:phyxminiproblems-problemphyxmini0147-hypotenusecontact, def:physics:phyx-mini-0147:phyxminiproblems-problemphyxmini0147-liquiddetectorprismsetup, def:physics:phyx-mini-0147:phyxminiproblems-problemphyxmini0147-matchesintendeddetectorlayout, def:physics:phyx-mini-0147:phyxminiproblems-problemphyxmini0147-hasintendedraygeometry, def:physics:phyx-mini-0147:phyxminiproblems-problemphyxmini0147-usesstandardairwaterindices, def:physics:phyx-mini-0147:phyxminiproblems-problemphyxmini0147-hasphysicalparameters, def:physics:phyx-mini-0147:phyxminiproblems-problemphyxmini0147-satisfiesprisminterfaceoptics, def:physics:phyx-mini-0147:phyxminiproblems-problemphyxmini0147-satisfiessensorresponse, def:physics:phyx-mini-0147:phyxminiproblems-problemphyxmini0147-allowableprismindices, def:physics:phyx-mini-0147:phyxminiproblems-problemphyxmini0147-recordeddatasetanswer, def:physics:phyx-mini-0147:phyxminiproblems-problemphyxmini0147-matchesroundedallowablerange}
The assigned autoformalize agent should translate this physics problem into Lean declarations in the covered file, with theorem and lemma proof bodies written as `by sorry`.
\end{theorem}
\begin{proof}
This is an autoformalization task, not a proof task. Produce statements that can later be proved by the physics prover without weakening the physical model.
\end{proof}

% --- Archon declaration topology begin ---
\section{Declaration topology}
\begin{definition}[Prism Face]
  \label{def:physics:phyx-mini-0147:phyxminiproblems-problemphyxmini0147-prismface}
  \lean{PhyXMiniProblems.ProblemPhyXMini0147.PrismFace}
  The three prism faces met by, or relevant to, the detector beam
\end{definition}
\begin{proof}
  This is the declared model definition of Prism Face.
\end{proof}
\begin{definition}[Prism Cross Section]
  \label{def:physics:phyx-mini-0147:phyxminiproblems-problemphyxmini0147-prismcrosssection}
  \lean{PhyXMiniProblems.ProblemPhyXMini0147.PrismCrossSection}
  Shape of the prism cross-section intended by the detector description
\end{definition}
\begin{proof}
  This is the declared model definition of Prism Cross Section.
\end{proof}
\begin{definition}[Hypotenuse Contact]
  \label{def:physics:phyx-mini-0147:phyxminiproblems-problemphyxmini0147-hypotenusecontact}
  \lean{PhyXMiniProblems.ProblemPhyXMini0147.HypotenuseContact}
  Whether the externally accessible hypotenuse is dry or covered by liquid
\end{definition}
\begin{proof}
  This is the declared model definition of Hypotenuse Contact.
\end{proof}
\begin{definition}[Sensor Signal]
  \label{def:physics:phyx-mini-0147:phyxminiproblems-problemphyxmini0147-sensorsignal}
  \lean{PhyXMiniProblems.ProblemPhyXMini0147.SensorSignal}
  The two qualitative light-sensor readouts described in the problem
\end{definition}
\begin{proof}
  This is the declared model definition of Sensor Signal.
\end{proof}
\begin{definition}[degrees]
  \label{def:physics:phyx-mini-0147:phyxminiproblems-problemphyxmini0147-degrees}
  \lean{PhyXMiniProblems.ProblemPhyXMini0147.degrees}
  Convert a real-valued degree readout into a physical angle
\end{definition}
\begin{proof}
  This is the declared model definition of degrees.
\end{proof}
\begin{definition}[Prism Detector Diagram]
  \label{def:physics:phyx-mini-0147:phyxminiproblems-problemphyxmini0147-prismdetectordiagram}
  \lean{PhyXMiniProblems.ProblemPhyXMini0147.PrismDetectorDiagram}
  \uses{def:physics:phyx-mini-0147:phyxminiproblems-problemphyxmini0147-prismface, def:physics:phyx-mini-0147:phyxminiproblems-problemphyxmini0147-prismcrosssection}
  Combinatorial geometry of the intended detector diagram. The fields retain the source-facing leg, the sensor-facing leg, and the reservoir-contact face without assigning an answer interval to the prism index
\end{definition}
\begin{proof}
  This is the declared model definition of Prism Detector Diagram.
\end{proof}
\begin{definition}[Liquid Detector Prism Setup]
  \label{def:physics:phyx-mini-0147:phyxminiproblems-problemphyxmini0147-liquiddetectorprismsetup}
  \lean{PhyXMiniProblems.ProblemPhyXMini0147.LiquidDetectorPrismSetup}
  \uses{def:physics:phyx-mini-0147:phyxminiproblems-problemphyxmini0147-prismface, def:physics:phyx-mini-0147:phyxminiproblems-problemphyxmini0147-hypotenusecontact, def:physics:phyx-mini-0147:phyxminiproblems-problemphyxmini0147-sensorsignal, def:physics:phyx-mini-0147:phyxminiproblems-problemphyxmini0147-prismdetectordiagram}
  Optical observables of the detector for a hypothetical positive prism index. The external refractive index depends on whether the hypotenuse is exposed to air or covered by reservoir liquid. Optical outcomes and the sensor readout remain primitive observations; governing laws relate them below
\end{definition}
\begin{proof}
  This is the declared model definition of Liquid Detector Prism Setup.
\end{proof}
\begin{definition}[Matches Intended Detector Layout]
  \label{def:physics:phyx-mini-0147:phyxminiproblems-problemphyxmini0147-matchesintendeddetectorlayout}
  \lean{PhyXMiniProblems.ProblemPhyXMini0147.MatchesIntendedDetectorLayout}
  \uses{def:physics:phyx-mini-0147:phyxminiproblems-problemphyxmini0147-liquiddetectorprismsetup}
  The intended component layout: a right-isosceles prism, normal entry through one leg, sensor exit through the other leg, and liquid contact at the hypotenuse. This is a model assumption, not a readout of the mismatched bitmap
\end{definition}
\begin{proof}
  This is the declared model definition of Matches Intended Detector Layout.
\end{proof}
\begin{definition}[Has Intended Ray Geometry]
  \label{def:physics:phyx-mini-0147:phyxminiproblems-problemphyxmini0147-hasintendedraygeometry}
  \lean{PhyXMiniProblems.ProblemPhyXMini0147.HasIntendedRayGeometry}
  \uses{def:physics:phyx-mini-0147:phyxminiproblems-problemphyxmini0147-degrees, def:physics:phyx-mini-0147:phyxminiproblems-problemphyxmini0147-liquiddetectorprismsetup}
  Normal source entry leaves the internal ray undeviated. In a right-isosceles cross-section it then meets the hypotenuse at 45° from the face normal
\end{definition}
\begin{proof}
  This is the declared model definition of Has Intended Ray Geometry.
\end{proof}
\begin{definition}[Uses Standard Air Water Indices]
  \label{def:physics:phyx-mini-0147:phyxminiproblems-problemphyxmini0147-usesstandardairwaterindices}
  \lean{PhyXMiniProblems.ProblemPhyXMini0147.UsesStandardAirWaterIndices}
  \uses{def:physics:phyx-mini-0147:phyxminiproblems-problemphyxmini0147-liquiddetectorprismsetup}
  The dimensionless textbook environmental indices used by the recorded answer: n\_air = 1.00 and n\_water = 1.33. These are calibration assumptions, not the requested prism-index bounds
\end{definition}
\begin{proof}
  This is the declared model definition of Uses Standard Air Water Indices.
\end{proof}
\begin{definition}[Has Physical Parameters]
  \label{def:physics:phyx-mini-0147:phyxminiproblems-problemphyxmini0147-hasphysicalparameters}
  \lean{PhyXMiniProblems.ProblemPhyXMini0147.HasPhysicalParameters}
  \uses{def:physics:phyx-mini-0147:phyxminiproblems-problemphyxmini0147-liquiddetectorprismsetup}
  Positivity, optical ordering, and the acute physical incidence branch
\end{definition}
\begin{proof}
  This is the declared model definition of Has Physical Parameters.
\end{proof}
\begin{definition}[Satisfies Prism Interface Optics]
  \label{def:physics:phyx-mini-0147:phyxminiproblems-problemphyxmini0147-satisfiesprisminterfaceoptics}
  \lean{PhyXMiniProblems.ProblemPhyXMini0147.SatisfiesPrismInterfaceOptics}
  \uses{def:physics:phyx-mini-0147:phyxminiproblems-problemphyxmini0147-hypotenusecontact, def:physics:phyx-mini-0147:phyxminiproblems-problemphyxmini0147-liquiddetectorprismsetup}
  Critical-angle optics at the prism hypotenuse, uniformly for every candidate positive prism index. Strict total internal reflection occurs when the external index is below n\_prism sin θ. Light genuinely enters the external medium when that quantity is strictly below the external index. Equality is the grazing critical case and is excluded from both operating states. Neither criterion contains the requested numerical interval
\end{definition}
\begin{proof}
  This is the declared model definition of Satisfies Prism Interface Optics.
\end{proof}
\begin{definition}[Satisfies Sensor Response]
  \label{def:physics:phyx-mini-0147:phyxminiproblems-problemphyxmini0147-satisfiessensorresponse}
  \lean{PhyXMiniProblems.ProblemPhyXMini0147.SatisfiesSensorResponse}
  \uses{def:physics:phyx-mini-0147:phyxminiproblems-problemphyxmini0147-liquiddetectorprismsetup}
  Response law of the idealized sensor. Reflected light gives the large signal; transmission into the exterior decreases it. This law does not assert which contact state produces either optical outcome
\end{definition}
\begin{proof}
  This is the declared model definition of Satisfies Sensor Response.
\end{proof}
\begin{definition}[Is Allowable Prism Index]
  \label{def:physics:phyx-mini-0147:phyxminiproblems-problemphyxmini0147-isallowableprismindex}
  \lean{PhyXMiniProblems.ProblemPhyXMini0147.IsAllowablePrismIndex}
  \uses{def:physics:phyx-mini-0147:phyxminiproblems-problemphyxmini0147-liquiddetectorprismsetup}
  A candidate prism index is allowable precisely when it produces the two required physical operating modes: TIR while dry and transmission while wet. This definition is behavioral; it does not unfold to the requested bounds
\end{definition}
\begin{proof}
  This is the declared model definition of Is Allowable Prism Index.
\end{proof}
\begin{definition}[allowable Prism Indices]
  \label{def:physics:phyx-mini-0147:phyxminiproblems-problemphyxmini0147-allowableprismindices}
  \lean{PhyXMiniProblems.ProblemPhyXMini0147.allowablePrismIndices}
  \uses{def:physics:phyx-mini-0147:phyxminiproblems-problemphyxmini0147-liquiddetectorprismsetup, def:physics:phyx-mini-0147:phyxminiproblems-problemphyxmini0147-isallowableprismindex}
  All dimensionless prism indices for which the detector operates correctly
\end{definition}
\begin{proof}
  This is the declared model definition of allowable Prism Indices.
\end{proof}
\begin{definition}[Answer Choice]
  \label{def:physics:phyx-mini-0147:phyxminiproblems-problemphyxmini0147-answerchoice}
  \lean{PhyXMiniProblems.ProblemPhyXMini0147.AnswerChoice}
  Labels of the four refractive-index ranges displayed in the problem
\end{definition}
\begin{proof}
  This is the declared model definition of Answer Choice.
\end{proof}
\begin{definition}[displayed Lower Index Bound]
  \label{def:physics:phyx-mini-0147:phyxminiproblems-problemphyxmini0147-displayedlowerindexbound}
  \lean{PhyXMiniProblems.ProblemPhyXMini0147.displayedLowerIndexBound}
  \uses{def:physics:phyx-mini-0147:phyxminiproblems-problemphyxmini0147-answerchoice}
  Printed lower endpoint of each dimensionless answer interval
\end{definition}
\begin{proof}
  This is the declared model definition of displayed Lower Index Bound.
\end{proof}
\begin{definition}[displayed Upper Index Bound]
  \label{def:physics:phyx-mini-0147:phyxminiproblems-problemphyxmini0147-displayedupperindexbound}
  \lean{PhyXMiniProblems.ProblemPhyXMini0147.displayedUpperIndexBound}
  \uses{def:physics:phyx-mini-0147:phyxminiproblems-problemphyxmini0147-answerchoice}
  Printed upper endpoint of each dimensionless answer interval
\end{definition}
\begin{proof}
  This is the declared model definition of displayed Upper Index Bound.
\end{proof}
\begin{definition}[displayed Prism Index Interval]
  \label{def:physics:phyx-mini-0147:phyxminiproblems-problemphyxmini0147-displayedprismindexinterval}
  \lean{PhyXMiniProblems.ProblemPhyXMini0147.displayedPrismIndexInterval}
  \uses{def:physics:phyx-mini-0147:phyxminiproblems-problemphyxmini0147-answerchoice, def:physics:phyx-mini-0147:phyxminiproblems-problemphyxmini0147-displayedlowerindexbound, def:physics:phyx-mini-0147:phyxminiproblems-problemphyxmini0147-displayedupperindexbound}
  The open prism-index interval printed beside an answer label
\end{definition}
\begin{proof}
  This is the declared model definition of displayed Prism Index Interval.
\end{proof}
\begin{definition}[recorded Dataset Answer]
  \label{def:physics:phyx-mini-0147:phyxminiproblems-problemphyxmini0147-recordeddatasetanswer}
  \lean{PhyXMiniProblems.ProblemPhyXMini0147.recordedDatasetAnswer}
  \uses{def:physics:phyx-mini-0147:phyxminiproblems-problemphyxmini0147-answerchoice}
  The answer label recorded by the source dataset
\end{definition}
\begin{proof}
  This is the declared model definition of recorded Dataset Answer.
\end{proof}
\begin{definition}[Matches Rounded Allowable Range]
  \label{def:physics:phyx-mini-0147:phyxminiproblems-problemphyxmini0147-matchesroundedallowablerange}
  \lean{PhyXMiniProblems.ProblemPhyXMini0147.MatchesRoundedAllowableRange}
  \uses{def:physics:phyx-mini-0147:phyxminiproblems-problemphyxmini0147-answerchoice, def:physics:phyx-mini-0147:phyxminiproblems-problemphyxmini0147-displayedlowerindexbound, def:physics:phyx-mini-0147:phyxminiproblems-problemphyxmini0147-displayedupperindexbound}
  A displayed two-decimal interval matches exact physical endpoints when each endpoint is within half a hundredth. This expresses numerical rounding and does not make any particular choice correct by definition
\end{definition}
\begin{proof}
  This is the declared model definition of Matches Rounded Allowable Range.
\end{proof}
\begin{lemma}[hypotenuse Incidence Sine eq sqrt Two Over Two]
  \label{lem:physics:phyx-mini-0147:phyxminiproblems-problemphyxmini0147-hypotenuseincidencesine-eq-sqrttwoovertwo}
  \lean{PhyXMiniProblems.ProblemPhyXMini0147.hypotenuseIncidenceSine_eq_sqrtTwoOverTwo}
  \uses{def:physics:phyx-mini-0147:phyxminiproblems-problemphyxmini0147-liquiddetectorprismsetup, def:physics:phyx-mini-0147:phyxminiproblems-problemphyxmini0147-hasintendedraygeometry}
  The intended 45° incidence geometry gives sin θ = √2 / 2. This is a derived trigonometric relation
\end{lemma}
\begin{proof}
  The conclusion follows from the governing relations and figure readouts named in the statement of hypotenuse Incidence Sine eq sqrt Two Over Two.
\end{proof}
\begin{lemma}[allowable Prism Indices eq exact Interval]
  \label{lem:physics:phyx-mini-0147:phyxminiproblems-problemphyxmini0147-allowableprismindices-eq-exactinterval}
  \lean{PhyXMiniProblems.ProblemPhyXMini0147.allowablePrismIndices_eq_exactInterval}
  \uses{def:physics:phyx-mini-0147:phyxminiproblems-problemphyxmini0147-liquiddetectorprismsetup, def:physics:phyx-mini-0147:phyxminiproblems-problemphyxmini0147-hasintendedraygeometry, def:physics:phyx-mini-0147:phyxminiproblems-problemphyxmini0147-usesstandardairwaterindices, def:physics:phyx-mini-0147:phyxminiproblems-problemphyxmini0147-satisfiesprisminterfaceoptics, def:physics:phyx-mini-0147:phyxminiproblems-problemphyxmini0147-allowableprismindices}
  The critical-angle and transmission laws characterize the exact, unrounded operating range. Air-side TIR gives √2 < n, while transmission into water gives n < 1.33 √2
\end{lemma}
\begin{proof}
  The conclusion follows from the governing relations and figure readouts named in the statement of allowable Prism Indices eq exact Interval.
\end{proof}
\begin{lemma}[large Sensor Signal iff no Liquid]
  \label{lem:physics:phyx-mini-0147:phyxminiproblems-problemphyxmini0147-largesensorsignal-iff-noliquid}
  \lean{PhyXMiniProblems.ProblemPhyXMini0147.largeSensorSignal_iff_noLiquid}
  \uses{def:physics:phyx-mini-0147:phyxminiproblems-problemphyxmini0147-hypotenusecontact, def:physics:phyx-mini-0147:phyxminiproblems-problemphyxmini0147-liquiddetectorprismsetup, def:physics:phyx-mini-0147:phyxminiproblems-problemphyxmini0147-satisfiessensorresponse, def:physics:phyx-mini-0147:phyxminiproblems-problemphyxmini0147-isallowableprismindex}
  For any allowable prism index, the sensor has a large signal exactly in the absence of liquid. Thus the qualitative inference stated in the scenario is derived from the operating modes and sensor law, not assumed
\end{lemma}
\begin{proof}
  The conclusion follows from the governing relations and figure readouts named in the statement of large Sensor Signal iff no Liquid.
\end{proof}
% --- Archon declaration topology end ---
% --- Archon physics formalization source end ---
